% =============================================================================
\chapter{Objetivos}
% =============================================================================

\section{Objetivo geral}

Projetar e avaliar uma arquitetura híbrida de Game AI na qual uma política de
aprendizado por reforço multiobjetivo condicionada por preferências seleciona,
em tempo de execução, métodos de decomposição HTN entre as alternativas
aplicáveis.

\section{Objetivos específicos}

\begin{enumerate}
\item definir uma interface formal entre HTN, codificação de preferências e MORL;
\item modelar restrições obrigatórias por pré-condições e regras de aplicabilidade HTN;
\item representar preferências comportamentais flexíveis por vetores dinâmicos de pesos;
\item desenvolver uma política MORL que selecione métodos HTN aplicáveis;
\item implementar um ambiente controlado de jogo com recompensas vetoriais e eventos reproduzíveis;
\item comparar a proposta com mecanismos tradicionais de seleção de métodos;
\item medir desempenho, adaptação, diversidade, validade, interpretabilidade e custo computacional;
\item avaliar generalização para vetores de preferência e condições ambientais não vistos;
\item analisar a contribuição de cada componente por estudos de ablação;
\item comparar perfis fixos, regras explícitas e codificadores relacionais para
inferência de $\mathbf{w}_t$, mantendo fixa a política MORL quando aplicável.
\end{enumerate}

% =============================================================================
\chapter{Metodologia}
% =============================================================================

\section{Revisão da literatura}

Será conduzida uma revisão estruturada sobre HTN, GOAP como fonte de
representações relacionais, MORL, políticas condicionadas por preferências,
aprendizado por reforço hierárquico, arquiteturas neuro-simbólicas, NPCs
adaptativos e Game AI explicável.
A revisão deverá posicionar explicitamente a contribuição no nível de seleção
de métodos HTN e identificar trabalhos que aprendem políticas de baixo nível,
métodos de planejamento ou preferências.

\section{Formalização}

A formalização deverá separar quatro elementos: validade semântica,
preferência comportamental, utilidade aprendida e execução de baixo nível.
Serão definidos o estado, as tarefas, os métodos, a função de aplicabilidade, o
vetor de recompensas, o vetor de preferências, a política condicionada e os
critérios de término de uma decisão hierárquica.

\section{Ambiente e dados}

Os dados serão gerados por simulação.
Cada episódio registrará estados, tarefas, métodos, recompensas, preferências,
resultados, consumo de recursos, dano, tempo e latência.
O ambiente deverá permitir cenários repetíveis, configuração de eventos estocásticos,
execução automatizada, sementes determinísticas e visualização de traços de decisão.

O domínio inicial deverá conter um agente principal, três a cinco tarefas
compostas, duas a cinco decomposições por tarefa e três a cinco objetivos.
Um cenário tático ou de sobrevivência é adequado por produzir compromissos
claros entre missão, segurança, aliados, recursos e tempo.

\section{Implementação HTN}

O HTN definirá a estrutura válida do comportamento.
O planner será adaptado para calcular $M_{\mathrm{valid}}(s_t,\tau_t)$ e permitir que um
seletor externo escolha a alternativa.
As restrições obrigatórias permanecerão exclusivamente na camada simbólica.

\section{Codificação de preferências}

Perfis fixos e regras explícitas converterão estado e contexto em um vetor
normalizado e servirão como baselines.
Exemplos incluem aumentar sobrevivência sob baixa saúde, proteção do aliado
quando seu estado se torna crítico e economia de recursos quando a munição
diminui.
Em seguida, será avaliado um codificador baseado em grafo relacional que
represente objetivos, condições, efeitos, recursos e ameaças; essa representação
é inspirada em GOAP, mas não executa um planejador GOAP. Um modelo profundo será
incluído somente se houver dados, orçamento e evidência experimental de que a
complexidade adicional é justificável.
Todos os codificadores serão avaliados sob a mesma interface $E(g_t,s_t,p,c_t)\mapsto\mathbf{w}_t$.

\section{Treinamento MORL}

A política receberá estado, tarefa composta, máscara de métodos válidos e vetor
de preferência.
A ação será a escolha de um método.
O treinamento amostrará estados iniciais, configurações de cenário, eventos e
vetores de preferência.
O orçamento de treinamento, os hiperparâmetros e as escalas de recompensa serão
controlados e reportados.

A atribuição de crédito será definida no nível da decisão hierárquica.
O retorno poderá acumular recompensas durante a execução das subtarefas produzidas
pelo método até a próxima escolha HTN, término da decomposição ou interrupção por
replanejamento.

\section{Baselines}

\begin{description}
\item[Baseline 1 --- primeiro aplicável:] seleciona o primeiro método válido na ordem do domínio.
\item[Baseline 2 --- prioridade estática:] utiliza prioridades manuais constantes.
\item[Baseline 3 --- utilidade com pesos fixos:] combina avaliações manuais por um único vetor fixo.
\item[Baseline 4 --- utilidade dinâmica baseada em regras:] usa regras explícitas para $\mathbf{w}_t$, mas valores manuais, sem aprendizado.
\item[Baseline 5 --- MORL com perfis ou regras:] usa a política MORL com um codificador simples de preferências.
\end{description}

Esses baselines separam o ganho produzido por preferências dinâmicas do ganho
produzido por valores aprendidos.
Codificadores relacionais e profundos serão comparados à configuração MORL equivalente,
sem alterar a máscara HTN.

\section{Estudos de ablação}

Serão avaliadas, quando tecnicamente viáveis, as seguintes configurações:

\begin{itemize}
\item MORL com preferências fixas;
\item MORL sem filtragem HTN, em ambiente seguro ou com penalidades para invalidez;
\item HTN com perfis fixos ou regras, sem MORL;
\item HTN e MORL sem mudanças dinâmicas de preferência;
\item arquitetura HTN--codificador de preferências--MORL;
\item codificador relacional versus regras explícitas sob a mesma política MORL.
\end{itemize}

\section{Procedimento experimental}

Cada arquitetura será executada sobre conjuntos equivalentes de cenários e
sementes.
Serão mantidos, tanto quanto possível, os mesmos horizontes, orçamentos
computacionais e protocolos de ajuste.
Os logs serão agregados por arquitetura e condição.
A análise incluirá estatística descritiva, intervalos de confiança, testes
adequados à distribuição dos dados e tamanhos de efeito.

Mudanças de contexto e preferência serão introduzidas em pontos controlados do
episódio para medir adaptação.
Vetores não usados diretamente no treinamento e cenários com combinações novas
serão reservados para avaliação de generalização.

% =============================================================================
\chapter{Métricas de Avaliação}
% =============================================================================

\section{Desempenho da tarefa}

\begin{itemize}
\item taxa de sucesso da missão e sobrevivência;
\item conclusão de objetivos e sobrevivência de aliados;
\item dano recebido e recursos consumidos;
\item duração do episódio e número de replanejamentos.
\end{itemize}

\section{Desempenho multiobjetivo}

As métricas serão selecionadas conforme as hipóteses do cenário e reportadas
com suas referências, em particular o protocolo de benchmarking MORL
\cite{felten2023toolkit}.

\begin{itemize}
\item retorno vetorial acumulado;
\item utilidade escalarizada sob diferentes preferências;
\item hipervolume e cobertura de soluções somente quando houver conjunto de soluções
e ponto de referência explicitamente definidos;
\item utilidade esperada sobre uma distribuição de preferências;
\item \textit{Maximum Utility Loss} (MUL), quando houver solução de referência;
\item \textit{Inverted Generational Distance} (IGD), apenas em domínios com
fronteira de Pareto de referência conhecida;
\item desempenho sob vetores de preferência não vistos.
\end{itemize}

Para utilidade linear, a utilidade esperada será reportada apenas quando a
distribuição sobre preferências estiver explicitamente definida; métricas de
hipervolume não substituem a avaliação sob preferências condicionadas.

\section{Adaptação e generalização}

\begin{itemize}
\item desempenho após mudanças contextuais;
\item latência de adaptação após alteração de preferências;
\item frequência e adequação das mudanças de método;
\item recuperação após eventos inesperados;
\item desempenho em cenários e combinações não observados no treinamento.
\end{itemize}

\section{Diversidade e coerência comportamental}

\begin{itemize}
\item distribuição dos métodos HTN selecionados;
\item número de estratégias distintas observadas;
\item variação entre personalidades e vetores de preferência;
\item consistência sob estados e preferências equivalentes;
\item taxa de oscilação ou troca prematura de estratégia.
\end{itemize}

\section{Validade semântica}

\begin{itemize}
\item taxa de seleção de método inválido;
\item violações de pré-condições;
\item tentativas de uso de recursos indisponíveis;
\item violações de sequência lógica de tarefas.
\end{itemize}

Para a arquitetura completa, espera-se taxa de seleção inválida igual a zero,
pois os métodos não aplicáveis são removidos antes da decisão MORL. Falhas de
decomposição, falhas de execução e replanejamentos devem ser contabilizados
separadamente.

\section{Desempenho computacional}

\begin{itemize}
\item latência total de decisão, tempo da fonte de preferências e tempo de inferência;
\item tempo de filtragem HTN;
\item tempo de decomposição do método selecionado;
\item taxa e custo de \emph{method fallback};
\item tempo de treinamento;
\item quantidade de métodos decompostos no caminho normal;
\item consumo de memória;
\item estabilidade da frequência de decisão em tempo real.
\end{itemize}

\section{Densidade de decisão}

Além de medir latência, a avaliação reportará a densidade de decisão: número de
consultas ao seletor por episódio, por segundo simulado e por ação primitiva.
Essa métrica testa se a abstração HTN reduz a frequência efetiva de escolhas sem
presumir, a priori, que seu espaço de decisão é globalmente menor.

\section{Rastreabilidade e auditabilidade}

Cada decisão será registrada com tarefa composta, métodos aplicáveis, métodos
rejeitados e pré-condições falhas, vetor de preferências, valores objetivos,
utilidade escalarizada e método escolhido.
A avaliação verificará completude dos traços e coerência entre prioridades e
decisão, complementada por inspeção qualitativa de episódios representativos.

